\documentclass[11pt]{article}
\usepackage[margin=1in]{geometry}
\usepackage{amsmath}
\usepackage[hidelinks]{hyperref}

\title{Kinematic Basis of the Linkage Simulator Prototype}
\author{Implementation-derived technical note}
\date{}

\begin{document}
\maketitle

\section{Scope and conventions}

The prototype is a deterministic planar kinematics demonstrator, not a dynamic or
biomechanical model. World coordinates use millimetre-like units, angles use
radians internally, and positive angles are counter-clockwise. A rigid link of
length $L$ has a centre pose $(\mathbf{p},\theta)$ and local endpoints
$(-L/2,0)$ and $(L/2,0)$. A local point $\mathbf{q}$ maps to world space as
\begin{equation}
  \mathbf{x} = \mathbf{p} + R(\theta)\mathbf{q}, \qquad
  R(\theta)=\begin{bmatrix}\cos\theta&-\sin\theta\\
  \sin\theta&\cos\theta\end{bmatrix}.
\end{equation}
The data model permits arbitrary collections of links and revolute joints. The
first solver strategy is deliberately narrower: one four-bar topology identified
by stable component IDs.
\section{Dorsal mounting geometry}

Let $\mathbf{G}$ be a point on the dorsal hand-adjacent ground plane, with
unit tangent $\mathbf{t}_g$ and dorsal normal $\mathbf{n}_g$. The editable
servo--ground offset $h$ places the servo axis at
\begin{equation}
  \mathbf{A}=\mathbf{G}+h\mathbf{n}_g.
\end{equation}
The fixed dorsal rail has editable length $L_g$ and servo-relative angle
$\alpha_g$. Its second endpoint, which grounds the proximal rocker, is
\begin{equation}
  \mathbf{D}=\mathbf{A}+L_g(\cos\alpha_g,\sin\alpha_g).
\end{equation}
Thus the rocker anchor is visibly structural rather than a floating ground
symbol. Repositioning the servo preserves the ground-plane normal reference;
editing $h$, $L_g$, or $\alpha_g$ deterministically rebuilds both ground
pivots before the mechanism solve.

\section{Analytic four-bar solve}

Let $\mathbf{A}$ be the servo pivot, $\mathbf{D}$ the rocker ground pivot, and
$r_1,r_2,r_3$ the crank length, crank-to-rocker joint distance along the
coupler, and rocker length. Both ground pivots are dorsal and proximal to the
D2 metacarpal; there is no distal ground pivot. Given servo angle $\phi$,
the crank endpoint is
\begin{equation}
  \mathbf{B}=\mathbf{A}+r_1(\cos\phi,\sin\phi).
\end{equation}
The floating output joint $\mathbf{C}$ must satisfy
\begin{equation}
  \|\mathbf{C}-\mathbf{B}\|=r_2,\qquad
  \|\mathbf{C}-\mathbf{D}\|=r_3.
\end{equation}
It is therefore obtained from the intersection of two circles. The geometry
routine explicitly distinguishes separate, contained, tangent, two-solution,
coincident, concentric, and negative-radius cases. For two solutions, the solver
chooses the candidate nearest the preceding valid $\mathbf{C}$. This local
continuity rule preserves the assembly branch during a servo sweep without
embedding a fixed ``upper'' or ``lower'' branch rule.

The output joint $\mathbf{C}$ is not grounded to the anatomy. Instead, a short
structural anchor leaves it at fixed coupler-relative offset $\delta$:
\begin{equation}
  \mathbf{E}=\mathbf{C}+L_a
  \begin{bmatrix}\cos(\theta_c+\delta)\\ \sin(\theta_c+\delta)\end{bmatrix},
  \qquad \theta_c=\operatorname{atan2}(C_y-B_y,C_x-B_x).
\end{equation}
Two revolute rectangular drivers continue from $\mathbf{E}$: one carries the
middle-phalanx contactor, and the second has an ungrounded right endpoint and
carries the distal contactor. This split accommodates two dorsal contacts
without introducing an anatomical distal ground attachment.

When no finite intersection exists, only the crank pose is updated. The previous
valid coupler and rocker poses remain available for rendering and the state is
marked invalid. No non-finite values enter the render loop.

\section{Simultaneous finger and contactor solve}

The index finger is a serial proximal--middle--distal chain. DIP flexion is
coupled to PIP flexion by $\gamma=k\beta$, with $k=0.62$. For contactor $j$,
let $\mathbf{F}_j(\alpha,\beta)$ be its point on the finger centreline and
$\mathbf{n}_{d,j}$ the segment normal directed toward positive world $Y$.
The required linkage centreline point is
\begin{equation}
  \mathbf{T}_j=\mathbf{F}_j+
  \left(\frac{w_{f,j}}{2}+\frac{w_{l,j}}{2}+c_j\right)\mathbf{n}_{d,j},
\end{equation}
where $c_j$ is band clearance. The unknown vector
$\mathbf{z}=(\alpha,\beta,\eta_m,\eta_d)$ contains MCP and PIP angles plus the
two passive driver angles. With $\mathbf{P}_m(\mathbf{z})$ and
$\mathbf{P}_d(\mathbf{z})$ denoting the two linkage contact points, the default
solve drives the four residuals
\begin{equation}
  \mathbf{r}(\mathbf{z})=
  \begin{bmatrix}
    \mathbf{P}_m-\mathbf{T}_m\\
    \mathbf{P}_d-\mathbf{T}_d
  \end{bmatrix}=\mathbf{0}.
\end{equation}
A bounded Newton iteration uses a finite-difference $4\times4$ Jacobian and the
previous valid pose as its initial state. This is a determinate positional IK;
no artificial joint stiffness or force model is introduced. Rendering places a
rectangular pad dorsal to each digit segment and a U-shaped band down to the
flexor side.

Every rectangular link is also tested against each phalanx capsule. For link
centreline segment $S_l$, finger centreline segment $S_f$, link width $w_l$,
finger width $w_f$, and clearance $\epsilon$, a pose is admissible only if
\begin{equation}
  \operatorname{dist}(S_l,S_f) \geq \frac{w_l+w_f}{2}+\epsilon.
\end{equation}
The contact solve includes each link half-width in its dorsal offset, so the
intended bars remain outside the finger capsules. Links are also tested against
the finite dorsal ground-plane segment and base rail. Within the rail's
longitudinal span, a signed half-plane test prevents a servo-driven link from
passing to the hand-side of the rail even when it does not intersect the rail.
Intended servo and rocker pivot coincidences are permitted; other rail crossings
or capsule overlaps are rejected. Rejection restores the last valid link, hand,
mount, and contact configuration.

\section{Inverse endpoint controls}

Each rectangular link exposes its positive-local-$x$ endpoint as a control
handle. Dragging the crank handle directly sets and clamps the servo angle.
For constrained four-bar, anchor, middle-driver, and distal-driver handles, the
prototype samples the servo ROM and selects the valid angle whose endpoint is
nearest the pointer target.
This bounded one-dimensional inverse solve preserves the same forward solver as
the source of truth. An unconstrained imported link instead rotates about its
left endpoint.

User-added links are never floating: creation requires an underlying link or
servo context and atomically creates a revolute joint at the new rectangle's
left endpoint. Passive parent--child attachments are propagated after the main
four-bar solve so their joint locations follow moving parent links.

\section{Update ordering and invariants}

Simulation time advances independently from drawing. The servo sweeps between
its configured limits, reflecting direction at each limit, before the analytic
mechanism and finger solves run. Rendering is read-only with respect to mechanism
state. The principal invariants are: link poses remain finite; solved endpoint
distances agree with link lengths up to floating-point tolerance; and branch
choice is based only on the previous valid configuration.

\section{Next mathematical milestone}

A general constraint graph should replace the four-bar strategy. A practical next
step is a graph traversal that resolves analytically solvable dyads and reports
degrees of freedom, disconnected components, singularities, and overconstraints.
Numerical residual minimization can then serve as a bounded fallback, with the
analytic result and prior frame used as its initial condition.

\end{document}
